\section{Design philosophy \& the longitudinal convention}\label{sec:philosophy}

Before a single type is introduced, the design fixes the small set of invariants that
everything downstream must respect. These are deliberately stated as rules rather than
as code, because their value is exactly that they constrain every later decision: the
data model (\cref{sec:datamodel}), the resolver (\cref{sec:resolver}), the
serialization format (\cref{sec:serialization}), and the migration (\cref{sec:migration})
are all forced to be consistent with them. Two things are settled here: the three
non-negotiable rules that govern \emph{how} placement is represented and computed, and
the single longitudinal sign convention that makes the geometry unambiguous.

\subsection{The three non-negotiable rules}

The central defect of the legacy link (\cref{sec:intro}) is that it conflates two ideas
that must be kept apart: the \emph{author's intent} (``this rim seats against that rim'')
and a \emph{derived numeric placement} (the resolved axial offset). The legacy
\code{Vector3} stored a derived quantity --- a center-of-mass-to-center-of-mass offset ---
and then asked the author to hand-compute it. The rules below exist to make that class of
error structurally impossible.

\begin{designrule}[The three non-negotiable rules]
\textbf{1. Placement is \emph{derived}, recomputed by exactly one resolver.}
Absolute coordinates are never stored, and the resolved pose is never stored. There is a
single function (\cref{sec:resolver}) that turns intent into geometry, and both the
simulator and the visualizer read its output. \\[2pt]
\textbf{2. The stored link is physical \emph{intent}} --- a feature-to-feature relationship
between a station on the parent and a station on the child --- never an opaque offset. The
center of mass falls out as a consequence of the geometry, not as an authored input. \\[2pt]
\textbf{3. The common case costs zero authored numbers.} A child attached with no explicit
link simply abuts its parent; the trivial nose\,$\rightarrow$\,body\,$\rightarrow$\,fins
stack requires no coordinates at all (\cref{sec:authoring}).
\end{designrule}

Rule~1 is not a new discipline imposed from outside; it is the discipline the existing
codebase already follows for every other derived quantity. The Barrowman aerodynamic
contribution is documented as a ``\emph{Pure function of geometry; NOT stored (no staleness,
no \texttt{-Werror} field)}'' (\src{model/parts/Part.h}{154}), and the
center-of-mass-versus-middle offset is explicitly marked ``\emph{NOT CURRENTLY CONSUMED}''
and recomputed rather than cached (\src{model/parts/Part.h}{310}). Composite mass, CM, and
inertia are themselves rebuilt lazily behind a dirty flag rather than persisted. Placement
is simply the last quantity to be brought under the same rule.

\begin{rejected}[Storing the resolved transform]
The obvious shortcut --- resolve once, store the resulting transform, and read it back ---
was considered and rejected. It reintroduces a staleness class that is the exact
\emph{inverse} of today's bug: edit a part's length and its stored transform is silently
wrong, instead of (as today) edit an offset and the geometry is silently wrong. Either
direction leaves two sources of truth that can disagree. The design keeps a single source
of truth --- intent --- and re-derives geometry on every read, accepting a recompute cost
that the caching split in \cref{sec:resolver} keeps off the per-step hot path.
\end{rejected}

\subsection{Settling the longitudinal convention}

A placement model needs an unambiguous axial sign convention: a part's local origin must
sit at a defined plane, and ``forward'' must mean one direction. The design fixes
\textbf{\zfwd} --- the local origin of every part lies on the longitudinal axis at its
\emph{aft} plane, so the part occupies $z \in [0, +\text{length}]$ moving forward toward the
nose tip.

This is not an arbitrary choice. It is the \emph{cheapest correct} choice, because both the
renderer and the existing design files already use it; choosing the opposite would impose a
sign-inverting rewrite on both with no offsetting benefit.

\begin{itemize}
  \item \textbf{The visualizer is already \zfwd{}.} \code{buildCone} places the cone base at
  $z_{\text{base}} = -L/2$ and the tip at $z_{\text{tip}} = +L/2$
  (\src{visualizer/RocketMesh.cpp}{93}), with the base-cap normal pointing $-z$ (``outward
  $=$ aft'', \src{visualizer/RocketMesh.cpp}{132}). \code{buildTube} agrees, naming its
  $+z$ end ``forward'' and its $-z$ end ``aft'' (\src{visualizer/RocketMesh.cpp}{156}). The
  sphere is built from a $+z$ pole to a $-z$ pole on the same axis
  (\src{visualizer/RocketMesh.cpp}{294}). Every primitive the renderer already emits treats
  $+z$ as forward.
  \item \textbf{Legacy \texttt{.qrd} files are already \zfwd{}.} In the
  \texttt{xl75\_multi} fixture, the body is placed at a negative $z$ (extending aft, toward
  $-z$) and the coupler at a positive $z$ relative to the body (forward, toward the nose).
  The files encode the same sense the renderer does.
\end{itemize}

\begin{rejected}[$+z = $ aft]
Adopting \zaft{} would require a whole-axis inversion of the renderer --- every cone
base-cap normal, every sphere pole, every primitive endpoint --- \emph{and} a sign-negating
migration of every legacy file. Under \zfwd{}, by contrast, wiring the renderer to the new
resolver is a pure \emph{consumer swap} with no sign inversion, and the file migration
(\cref{sec:migration}) is a re-expression of the same numbers with no axial flip. The
alternative pays a large, error-prone cost to arrive at an equivalent model.
\end{rejected}

\subsection{The per-part-type local frame}

With \zfwd{} fixed, each concrete part type's local frame is fully determined. The aft plane
is always the origin; the part extends forward to its length.

\begin{itemize}
  \item \textbf{\code{ConicalNoseCone}.} The aft plane (the base rim) is at $z = 0$ with
  $r = \text{baseRadius}$; the fore point (the tip) is at $z = +\text{length}$ with $r = 0$.
  The outer radius tapers linearly,
  $\text{radiusOuterAt}(z) = \text{baseRadius}\,(1 - z/\text{length})$.
  \item \textbf{\code{BodyTube}, \code{FinSet}, \code{HollowSphere}.} The aft plane is at
  $z = 0$ and the fore plane at $z = +\text{length}$.
\end{itemize}

\tikzfig{convention-zframe}{The settled \zfwd{} convention. Each part's local origin sits on
the longitudinal axis at its aft plane, occupying $z \in [0, +\text{length}]$. The single
frame collision --- the cone's internal CM frame, whose origin is at mid-length and whose
$+z$ points toward the base --- is shown opposed to the chosen forward direction and is
reconciled by an explicit adapter rather than by altering the cone's tensor
math.}{fig:convention}

\Cref{fig:convention} also depicts the one place where the chosen convention does not match a
pre-existing frame inside the code, which the next subsection treats as the single most
correctness-critical detail of the whole design.

\subsection{The one frame collision: the cone's internal CM frame}

The cone's inertia tensor is centroidal, and the helper that reports where that centroid sits
--- \code{coneCmOffset} (\src{model/parts/ConicalNoseCone.cpp}{49}) --- works in a frame that
is \emph{opposite} to the convention just settled. It places the cone in $[-L/2, +L/2]$ with
the base at $+L/2$, so its centroid is computed as
$z = +L/2 - \bar{h}$, with $\bar{h} = L/4$ for a solid cone and $L/3$ for a thin shell:

\begin{listing}[htbp]
\begin{minted}{cpp}
// The tensor is centroidal; report where that CM sits relative to the component middle so
// the assembly's addChildPart position is CM-to-CM. With the cone in [-L/2, +L/2] (base at
// +L/2), the centroid is hbar forward of the base: L/4 (solid), L/3 (shell) => z = +L/2 - hbar.
Vector3 ConicalNoseCone::coneCmOffset(double L, bool solid)
{
   const double hbar = solid ? (L / 4.0) : (L / 3.0);
   return Vector3{0.0, 0.0, L / 2.0 - hbar};
}
\end{minted}
\caption{The cone's internal CM frame (\srccap{model/parts/ConicalNoseCone.cpp}{49}):
mid-length origin, base at $+z$ --- the reverse of the chosen \zfwd{} placement frame, whose
origin is the aft (base) plane with the tip at $+z$.}\label{lst:conecmoffset}
\end{listing}

This frame (\cref{lst:conecmoffset}) is the only place in the codebase whose $+z$ disagrees
with the chosen forward direction. The tempting fix --- flip the cone's tensor and CM math so
its $+z$ matches --- is precisely the wrong move: that math is correct, is exercised by the
existing inertia tests, and is unrelated to placement. Inverting it to satisfy a placement
convention would risk a correctness regression in the physics core in order to avoid a
two-line coordinate transform.

\begin{keyinsight}[The single most correctness-critical detail]
The cone's internal CM frame (mid-length origin, base at $+z$) is the reverse of the chosen
\zfwd{} placement frame (aft-plane origin, tip at $+z$). The design reconciles the two with
\emph{one explicitly specified, unit-tested adapter} --- not by flipping the cone's tested
tensor math. In the placement frame the cone's CM lies at $\text{cmZ}_{\text{fwd}} = \bar{h}$,
i.e.\ $L/4$ forward of the base for a solid cone and $L/3$ for a shell, equivalently
$L - \bar{h}$ aft of the tip ($3L/4$ for a solid cone). This adapter (developed in
\cref{sec:datamodel}) is the \emph{only} consumer of \code{getCenterMassOffset()}
(\src{model/parts/Part.h}{96}); it lives behind the new \code{cmStationLocal()} virtual, which
the cone is the only part type to override for a frame flip. For a part whose CM is at mid-length
$\text{cmZ}_{\text{fwd}} = L/2$. Getting this one transform right is what makes the cone's CM
a correctly derived consequence of the geometry rather than a quietly inverted offset.
\end{keyinsight}

Isolating the collision to a single adapter is what allows the rest of the design to treat
placement as pure geometry. No implicit cross-frame vector addition survives anywhere else:
the resolver works entirely in the \zfwd{} placement frame, and the only translation between
frames happens at one named, tested boundary.
